\documentclass[11pt,a4paper]{article}
\usepackage{turkos}

\begin{document}
\phaseheader{13}{File Systems: VFS, initrd and TurkFS}{Turn a pile of numbered blocks into named files and directories, behind one interface that also covers devices and a read-only boot image}{83}{90}

\ataglance{3--4 weeks}{5}{Phase 12: block devices and the buffer cache; Phase 11: processes, \code{fork}/\code{exec}; Phase 9: mutexes}{A
virtual file system layer (in-memory inodes, open-file objects, per-process descriptor tables, path lookup, mount
points), a read-only root file system from a tar archive, \code{/dev} with \code{tty}, \code{null}, \code{zero} and
disks, and TurkFS: a writable on-disk file system in the MINIX V3 format that Linux's \code{fsck.minix} accepts as
clean. System calls: \code{open}, \code{close}, \code{read}, \code{write}, \code{lseek}, \code{dup}, \code{dup2},
\code{stat}, \code{fstat}, \code{getdents}, \code{mkdir}, \code{unlink}, \code{chdir}, \code{getcwd}.}

\section{Why this phase}

A disk is a very long array of blocks. Users want named files in a tree of directories that they can open, read,
grow and delete, and that survive crashes. The \textbf{file system} bridges the two, and it is the largest subsystem
you will write. It uses nearly everything from before: the buffer cache and block devices (Phase 12), processes and
\code{fork} semantics (Phase 11), user-pointer checks (Phase 10), and mutexes, because file-system operations sleep on
disk I/O while holding locks (Phase 9).

The phase has two halves. The \textbf{VFS} (virtual file system) is the file-system-independent layer: file
descriptors, open files, path names, mount points. It is what makes ``everything is a file'' possible in UNIX, since
a terminal, a disk and a regular file all look the same to \code{read}. Below it you plug in concrete file systems:
first a trivial read-only one for booting, then \textbf{TurkFS}, a real on-disk format with inodes, bitmaps and
indirect blocks. TurkFS uses the MINIX V3 layout from \OSDI, so the book's figures describe your disk exactly, and
Linux can check your images and even mount them.

\section{Concepts you need}

\subsection{Files, directories and descriptors}

\OSTEP chapter 39 describes the interface. A \textbf{file} is an array of bytes with a low-level name, its
\textbf{inode number}. A \textbf{directory} is a file whose contents are a list of (name, inode number) pairs, and
directories nest into a tree. A process doesn't use inodes directly; \code{open} returns a small integer, the
\textbf{file descriptor}, which indexes the process's descriptor table. Three levels of objects are involved
(Figure~\ref{fig:vfs}), and keeping them apart is the key to getting \code{fork} and \code{dup} right. The descriptor
table is per process. An \textbf{open file} records the current offset and access mode and can be shared, so after
\code{fork} the parent and child share one offset. An \textbf{in-memory inode} exists once per file no matter how many
times it is open.

\begin{figure}[htbp]
\centering
\begin{tikzpicture}[b/.style={draw=tkNavy!70, fill=tkLight, rounded corners=2pt, font=\sffamily\scriptsize, align=center, minimum height=0.9cm},
  every node/.style={font=\sffamily\scriptsize}]
  \node[b, fill=tkGreen!10, text width=2.2cm] (pa) at (0,1.0) {process A\\fd table: 0 1 2 \textbf{3}};
  \node[b, fill=tkGreen!10, text width=2.2cm] (pb) at (0,-1.0) {process B (child of A)\\fd table: 0 1 2 \textbf{3}};
  \node[b, fill=tkAmber!12, text width=2.3cm] (f) at (3.9,0) {\code{struct file}\\offset 4096, O\_RDWR\\refcnt 2};
  \node[b, fill=tkBlue!10, text width=2.4cm] (i) at (7.6,0) {\code{struct inode}\\ino 57, size 9000\\refcnt 1, ops = TurkFS};
  \node[b, fill=tkRed!8, text width=2.2cm] (sb) at (11.2,0.9) {mounted TurkFS\\superblock, bitmaps};
  \node[b, fill=tkGray!15, text width=2.2cm] (bc) at (11.2,-0.9) {buffer cache and block device \code{hdb}};
  \draw[tkarrow] (pa.east) -- (f); \draw[tkarrow] (pb.east) -- (f);
  \draw[tkarrow] (f) -- (i); \draw[tkarrow] (i) -- (sb); \draw[tkarrow] (i) -- (bc);
\end{tikzpicture}
\caption{After \code{fork}, descriptor 3 in both processes refers to the same open file, and so to the same offset.}
\label{fig:vfs}
\end{figure}

\subsection{How an on-disk file system is laid out}

\OSTEP chapter 40 builds a very simple file system, \emph{vsfs}, and TurkFS is its real-world sibling. The disk is
divided into regions: a \textbf{superblock} describing the whole file system, an \textbf{inode bitmap} and a
\textbf{zone bitmap} marking which inodes and data blocks are in use, the \textbf{inode table}, and the \textbf{data}
region (Figure~\ref{fig:layout}). MINIX calls data blocks \emph{zones}; with the settings you will use, a zone is
exactly one 1\,KiB block (\OSDI\,\S5.6.2).

\begin{figure}[htbp]
\centering
\begin{tikzpicture}[r/.style={draw=tkNavy!70, minimum height=0.9cm, font=\sffamily\scriptsize, align=center}]
  \node[r, fill=tkGray!15, minimum width=1.1cm, anchor=west] (b0) at (0,0) {boot};
  \node[r, fill=tkRed!10, minimum width=1.3cm, right=0pt of b0] (b1) {super-\\block};
  \node[r, fill=tkAmber!14, minimum width=1.3cm, right=0pt of b1] (b2) {inode\\bitmap};
  \node[r, fill=tkAmber!14, minimum width=1.3cm, right=0pt of b2] (b3) {zone\\bitmap};
  \node[r, fill=tkBlue!10, minimum width=3.3cm, right=0pt of b3] (b4) {inode table\\(64-byte inodes, 16 per block)};
  \node[r, fill=tkGreen!10, minimum width=5.0cm, right=0pt of b4] (b5) {data zones\\(file contents, directories, indirect blocks)};
  \foreach \n/\t in {b0/0, b1/1, b2/2, b3/3, b4/4--174, b5/175--8191}
    \node[font=\ttfamily\scriptsize, text=tkGray, above=1pt] at (\n.north) {\t};
\end{tikzpicture}
\caption{TurkFS on an 8\,MiB image, exactly as \code{mkfs.minix -3} lays it out: 2736 inodes, first data zone 175. Numbers are 1\,KiB block numbers.}
\label{fig:layout}
\end{figure}

\begin{center}
\small\renewcommand{\arraystretch}{1.2}
\begin{tabularx}{\linewidth}{@{}l >{\raggedright\arraybackslash}X@{}}
\toprule
\tkth{Structure} & \tkth{Format (all little-endian)} \\
\midrule
superblock (at byte 1024) & \code{u32 s_ninodes; u16 pad; u16 s_imap_blocks; u16 s_zmap_blocks; u16 s_firstdatazone; u16 s_log_zone_size; u16 pad; u32 s_max_size; u32 s_zones; u16 s_magic (0x4D5A); u16 pad; u16 s_blocksize (1024); u8 s_disk_version} \\
inode (64 bytes) & \code{u16 i_mode; u16 i_nlinks; u16 i_uid; u16 i_gid; u32 i_size; u32 i_atime, i_mtime, i_ctime; u32 i_zone[10]} \\
directory entry (64 bytes) & \code{u32 inode; char name[60]} (name padded with zeros; inode 0 means an empty slot) \\
bitmaps & bit $n$ of the inode map is inode $n$; bit $n$ of the zone map is zone $n + \code{s_firstdatazone} - 1$; bit 0 of both is reserved and always set \\
\bottomrule
\end{tabularx}
\end{center}

\subsection{Finding a file's data}

An inode has ten zone pointers (Figure~\ref{fig:inode}). The first seven point directly at data zones, so small files
need nothing else. Pointer 7 points to a \textbf{single indirect} block holding 256 more zone numbers, and pointer 8 to
a \textbf{double indirect} block holding 256 pointers to blocks of 256 zone numbers each. The tenth pointer (triple
indirect) is reserved and TurkFS doesn't use it. \OSTEP\,\S40.3 and \OSDI\,\S5.6.4 explain the scheme, and \MOS\,\S4.3.2
compares it with the alternatives: contiguous allocation, linked lists and FAT.

\begin{figure}[htbp]
\centering
\begin{tikzpicture}[zptr/.style={draw=tkNavy!70, fill=tkBlue!8, minimum width=2.4cm, minimum height=0.42cm, font=\ttfamily\scriptsize},
  d/.style={draw=tkNavy!70, fill=tkGreen!12, minimum width=1.3cm, minimum height=0.42cm, font=\sffamily\scriptsize},
  ind/.style={draw=tkNavy!70, fill=tkAmber!14, minimum width=1.9cm, minimum height=0.85cm, font=\sffamily\scriptsize, align=center}]
  \node[font=\sffamily\scriptsize\bfseries] at (0,0.55) {inode};
  \node[zptr] (z0) at (0,0) {i\_zone[0]};
  \node[zptr, below=0pt of z0] (z1) {i\_zone[1..5]};
  \node[zptr, below=0pt of z1] (z6) {i\_zone[6]};
  \node[zptr, below=0pt of z6] (z7) {i\_zone[7]};
  \node[zptr, below=0pt of z7] (z8) {i\_zone[8]};
  \node[zptr, below=0pt of z8, fill=tkGray!12] (z9) {i\_zone[9] (unused)};
  \node[d] (d0) at (3.3,0) {data};
  \node[d] (d6) at (3.3,-0.84) {data};
  \node[ind] (s) at (3.6,-2.0) {single indirect\\256 zone numbers};
  \node[d] (sd) at (6.4,-2.0) {256 data zones};
  \node[ind] (dd) at (3.6,-3.3) {double indirect\\256 pointers};
  \node[ind] (dd2) at (6.6,-3.3) {256 blocks of\\256 zone numbers};
  \node[d, minimum width=2.1cm] (ddd) at (9.6,-3.3) {65\,536 data zones};
  \draw[tkarrow] (z0.east) -- (d0.west); \draw[tkarrow] (z6.east) -- (d6.west);
  \draw[tkarrow] (z7.east) -- (s.west); \draw[tkarrow] (s) -- (sd);
  \draw[tkarrow] (z8.east) -- (dd.west); \draw[tkarrow] (dd) -- (dd2); \draw[tkarrow] (dd2) -- (ddd);
\end{tikzpicture}
\caption{Zone pointers in a TurkFS inode, with 1\,KiB zones and 4-byte zone numbers.}
\label{fig:inode}
\end{figure}

\subsection{Path lookup, and what it costs}

To open \code{/bin/hello}, the kernel starts at the root inode (inode 1), reads the root directory's data and
searches it for \code{bin}, reads inode \code{bin}, reads its data and searches for \code{hello}, then reads inode
\code{hello}. \OSTEP\,\S40.4 draws this as a timeline of disk reads, and it is the best way to understand why caching
matters: without the buffer cache, every \code{open} costs several disk reads just to walk the path.

\subsection{Crashes}

Creating a file changes several blocks: the inode bitmap, the inode, the directory, maybe the zone bitmap and a data
block. If the machine stops half-way, the file system is left inconsistent: an inode marked used that no directory
points to, or a zone used by two files. \OSTEP chapter 42 presents the two classic answers. \textbf{fsck} scans the
whole disk after a crash and repairs it, which is simple but slow. \textbf{Journaling} writes a description of the
change to a log first and replays it after a crash. TurkFS relies on \code{fsck.minix} on the host for now and orders
its writes carefully (allocate, then initialise, then link), so a crash leaves at worst leaked blocks instead of
corruption. A journal is a stretch goal.

\section{Reading plan}

\begin{readingplan}
\must & \LOB & Ch.~12 \emph{File Systems}: a read-only file system, inodes, a virtual file system & 69--70 & The route of this phase in two pages. \\
\must & \OSTEP & Ch.~39 \emph{Interlude: Files and Directories} & 457--476 & The interface: descriptors, offsets, \code{fork} and \code{dup} sharing, links. \\
\must & \OSTEP & Ch.~40 \emph{File System Implementation} (vsfs, inodes, multi-level index, the access-path timeline, caching) & 477--494 & The design of TurkFS in miniature. \\
\must & \OSTEP & Ch.~42 \emph{Crash Consistency: FSCK and Journaling} & 507--526 & Why write order matters. \\
\must & \OSDI & \S5.6 \emph{Overview of the MINIX 3 File System}: messages, layout, bitmaps, i-nodes, block cache, directories and paths, descriptors, locking, pipes, the READ walkthrough & 548--566 & This is your on-disk format, explained by its designers. \\
\should & \OSDI & \S5.7 \emph{Implementation of the MINIX 3 File System} & 566--606 & Working code for every operation in Tasks 13.5--13.8. \\
\should & \MOS & \S4.1--4.3 \emph{Files}, \emph{Directories}, \emph{File-System Implementation} (incl.\ journaling and VFS) & 265--299 & Allocation methods compared; VFS in \S4.3.7. \\
\should & \MOS & \S4.5 \emph{Example File Systems} (MS-DOS FAT, UNIX V7) & 320--331 & Two real designs next to yours. \\
\should & \OSC & Ch.~13 \emph{File-System Interface}; \S14.1--14.5; \S15.5 \emph{Virtual File Systems} & 529--582, 603--605 & A thorough second treatment. \\
\could & \OSTEP & Ch.~41 \emph{Locality and the Fast File System}; Ch.~43 \emph{Log-structured File Systems} & 495--506, 527--542 & How layout decisions affect speed. \\
\could & \MOS & \S4.4 \emph{File-System Management and Optimization} & 299--320 & Block size, free-space tracking, consistency checks, caching. \\
\end{readingplan}

\begin{minixbox}
The MINIX file system server in \texttt{servers/fs/} implements the exact format you are writing. Read it alongside
your code: the superblock and bitmap allocation (\code{alloc_bit}, \code{free_bit}) in \texttt{super.c (23300)}; the
in-memory inode table (\code{get_inode}, \code{put_inode}, \code{alloc_inode}) in \texttt{inode.c (22900)}; path
lookup (\code{eat_path}, \code{advance}, \code{search_dir}) in \texttt{path.c (26300)}; file position to zone
mapping (\code{read_map}) in \texttt{read.c (25000)}; \code{open} and \code{mkdir} in \texttt{open.c (24500)};
descriptors in \texttt{filedes.c (23700)}. \OSDI\,\S5.6.10 traces one \code{read} call through all of it.
\end{minixbox}

\section{Build it}

\task{The VFS objects}
Define the three core structures. \code{struct inode} is the in-memory inode, file-system independent, with a table
of operations that each file system fills in. \code{struct file} is an open file. A mounted file system is a
\code{struct superblock} with its root inode. Keep in-memory inodes in a table keyed by (superblock, inode number), so
each file has exactly one, with a reference count.

\begin{lst}{c}{include/vfs.h}
struct inode_ops {
    int     (*lookup)(struct inode *dir, const char *name, struct inode **out);
    int     (*create)(struct inode *dir, const char *name, uint16_t mode, struct inode **out);
    int     (*mkdir)(struct inode *dir, const char *name, uint16_t mode);
    int     (*unlink)(struct inode *dir, const char *name);
    int     (*readdir)(struct inode *dir, uint32_t index, struct dirent *out);
    int32_t (*read)(struct inode *ino, void *buf, uint32_t n, uint32_t off);
    int32_t (*write)(struct inode *ino, const void *buf, uint32_t n, uint32_t off);
    int     (*truncate)(struct inode *ino, uint32_t size);
};

struct inode {
    uint32_t                ino;        /* number inside its file system        */
    uint16_t                mode;       /* type bits (dir, file, device) + rwx  */
    uint32_t                size;
    int                     refcnt;     /* in-memory references                 */
    struct superblock      *sb;
    const struct inode_ops *ops;
    struct mutex            lock;       /* held during operations that may sleep */
    struct inode           *mounted;    /* root of a file system mounted here   */
    void                   *priv;       /* file-system specific data            */
};

struct file {
    struct inode *inode;
    uint32_t      offset;
    int           flags;                /* O_RDONLY, O_WRONLY, O_RDWR, O_APPEND */
    int           refcnt;               /* descriptors pointing here            */
};
\end{lst}

\task{Descriptors and the basic calls}
Give each process a table of 32 \code{struct file} pointers. \code{open} resolves the path, allocates a
\code{struct file} and returns the lowest free descriptor. \code{read} and \code{write} copy through a kernel buffer
with \code{copy_to_user}/\code{copy_from_user} and advance the offset. \code{close} drops a reference. \code{dup} and
\code{dup2} copy a pointer and increment the count. \code{fork} copies the whole table, incrementing every count;
\code{exec} keeps it; \code{exit} closes everything. Keep using the Linux i386 numbers.

\begin{center}
\small\renewcommand{\arraystretch}{1.2}
\begin{tabular}{@{}l l l l l l@{}}
\toprule
\tkth{Call} & \tkth{No.} & \tkth{Call} & \tkth{No.} & \tkth{Call} & \tkth{No.} \\
\midrule
\code{open} & 5 & \code{lseek} & 19 & \code{stat} & 106 \\
\code{close} & 6 & \code{mkdir} & 39 & \code{fstat} & 108 \\
\code{unlink} & 10 & \code{dup} & 41 & \code{getdents} & 141 \\
\code{chdir} & 12 & \code{dup2} & 63 & \code{getcwd} & 183 \\
\bottomrule
\end{tabular}
\end{center}

\task{Devices as files}
Write a tiny \code{devfs} and mount it at \code{/dev}. Its inodes carry device operations instead of disk data:
\code{/dev/tty} reads a line from the keyboard and writes to the console; \code{/dev/null} discards writes and returns
end-of-file; \code{/dev/zero} returns zeros; \code{/dev/hda} and \code{/dev/hdb} read and write their block devices. When
the kernel starts \code{init}, it opens \code{/dev/tty} three times, so \code{init} and everything it starts inherit
descriptors 0, 1 and 2. \code{sys_write} to descriptor 1 now goes through the VFS like any other file.

\task{A read-only root from a tar archive}
The simplest file system to parse is a tar archive in the POSIX \emph{ustar} format, so use one as the initial root.
Each file is a 512-byte header followed by its contents padded to a multiple of 512 bytes, and the archive ends with
two zero blocks. The header holds the name (100 bytes at offset 0, plus a prefix of 155 bytes at offset 345), the size
as an octal ASCII string (12 bytes at offset 124), the type (one byte at 156: \code{'0'} for a file, \code{'5'} for a
directory) and the magic \code{"ustar"} at offset 257. Build the archive from a \code{rootfs/} directory in the
Makefile, load it as a boot module, and mount it at \code{/}. Now \code{execve} can take real paths.

\begin{lst}{make}{Makefile (root archive)}
$(BUILD)/initrd.tar: $(USER_PROGS) $(shell find rootfs -type f)
	mkdir -p rootfs/bin && cp $(USER_PROGS) rootfs/bin/
	tar --format=ustar -cf $@ -C rootfs .
\end{lst}
\verify{\code{execve("/bin/hello", ...)} loads the program from the archive, and \code{cat /etc/motd} (from the
monitor for now) prints a file you put in \code{rootfs/etc/}.}

\task{Path lookup}
Write \code{namei(path, &inode)}. Start at the root for absolute paths and at the process's current directory
otherwise, then take one component at a time: skip empty components (\code{//}), stay put on \code{.}, and call the
directory's \code{lookup} for everything else, including \code{..}. When the result has a file system mounted on it,
continue at that file system's root; when \code{..} is taken at the root of a mounted file system, go to the mount
point's parent. Also write \code{nameiparent} (the directory and the last component), which \code{create}, \code{mkdir}
and \code{unlink} need. Add \code{chdir} and \code{getcwd}.
\verify{\code{/bin/../etc/./motd}, \code{//etc//motd} and a relative \code{etc/motd} from \code{/} all resolve to the
same inode; \code{/dev/..} is \code{/}.}

\task{TurkFS, reading}
Create a file-system image on the host and attach it to QEMU as the primary slave (\code{hdb}); your ATA driver needs
to select the slave with \code{0xF0} in the drive register. Mount it at \code{/mnt}. Read the superblock through the
buffer cache and check the magic and block size; read inodes by number (inode $n$ is in block
$2 + \text{imap} + \text{zmap} + (n-1)/16$, at offset $((n-1) \bmod 16) \times 64$); map a file offset to a zone through the direct,
single and double indirect pointers; search directories entry by entry.

\begin{lst}{sh}{terminal (host)}
dd if=/dev/zero of=build/turkfs.img bs=1024 count=32768     # 32 MiB
mkfs.minix -3 build/turkfs.img                               # MINIX V3 layout, 1 KiB blocks
sudo mount -o loop build/turkfs.img /mnt                     # optional: put files on it
sudo cp -r rootfs/* /mnt && sudo umount /mnt
qemu-system-i386 ... -drive file=build/turkfs.img,format=raw,if=ide,index=1,media=disk
\end{lst}
\verify{\code{ls /mnt} (with the monitor, or a test program using \code{getdents}) lists what you copied from Linux,
and reading a file larger than 263\,KiB gives the right contents, which exercises the double indirect block.}

\task{TurkFS, writing}
Now allocation. Find a free inode by scanning the inode bitmap; find a free zone by scanning the zone bitmap and
remember the offset: bit $n$ is zone $n + \code{s_firstdatazone} - 1$. Implement \code{create} (allocate and
initialise an inode, then add a directory entry), \code{write} (allocate zones and indirect blocks as the file grows,
zeroing new blocks), \code{mkdir} (a new directory containing \code{.} and \code{..}, with the parent's link count
incremented), \code{unlink} (remove the entry, decrement the link count, and free the inode and its zones when the
count reaches zero and nobody has the file open), and \code{truncate}. Do every change through the buffer cache,
marking buffers dirty, and order the writes so a crash leaves at worst leaked space: allocate, initialise, and only then
link. Add a \code{sync} system call and flush all file systems at shutdown.
\verify{After Turk-OS creates, writes and deletes files, shut it down cleanly and run
\code{fsck.minix -f build/turkfs.img} on the host: it reports no errors and exits with status 0.}

\task{stat, fstat and getdents}
Define \code{struct stat} (device, inode, mode, link count, size, times) and \code{struct dirent} (inode number, name),
and implement the three calls through the inode operations. These are what \code{ls} will use in Phase 14.

\task{Tests}
Add these to \code{make test}, running on a fresh image each time; after the run, \code{make test} also runs
\code{fsck.minix} on the image.

\begin{center}
\small\renewcommand{\arraystretch}{1.25}
\begin{tabularx}{\linewidth}{@{}l X@{}}
\toprule
\tkth{Test} & \tkth{What must hold} \\
\midrule
\code{many_files} & create 200 files with random contents, read them back, delete them; free inode and zone counts return to their starting values \\
\code{big_file} & write and verify a 2\,MiB file (well into the double indirect range), then truncate it to 0 \\
\code{dirs} & nested \code{mkdir}, removing directories, link counts of parents \\
\code{paths} & \code{.}, \code{..}, \code{//}, relative paths, crossing into and out of \code{/mnt} and \code{/dev} \\
\code{fd_sharing} & after \code{fork}, parent and child share an offset; after \code{dup2}, two descriptors share one; closing one leaves the other working \\
\code{readonly_root} & writing to the tar root fails with \code{-EROFS} \\
\bottomrule
\end{tabularx}
\end{center}

\section{Testing and debugging}

When the on-disk format goes wrong, compare with what Linux sees. Mount the image read-only on the host
(\code{sudo mount -o loop,ro}) and look around, run \code{fsck.minix -fsv} for a detailed report, and dump raw blocks
with \code{hexdump -C -s $((BLOCK*1024)) -n 1024 build/turkfs.img}.

\begin{pitfalls}
\code{fsck.minix} reports zones marked in use but not used, or the reverse & The zone bitmap offset is wrong & Bit $n$ is zone $n + \code{s_firstdatazone} - 1$, not zone $n$. \\
Directory listing shows garbage names or skips entries & Inode 0 slots not skipped, or names treated as NUL-terminated when they are exactly 60 bytes & Skip entries with inode 0; copy at most 60 bytes and terminate yourself. \\
Wrong inode read & Inode numbers start at 1 & Use $n-1$ when computing the block and offset. \\
Writes vanish after reboot & Buffers not marked dirty, or no \code{sync} at shutdown & \code{bdirty} after every change; flush in the shutdown path. \\
After \code{fork}, the two processes read the file independently & The descriptor table was deep-copied into new \code{struct file} objects & Copy pointers and increment counts; they must share the offset. \\
Deadlock while creating a file & A directory inode lock and a buffer lock taken in different orders in different paths & Add the file-system locks to the lock order from Phase 9: inode, then buffer. \\
\end{pitfalls}

\section{Milestone}

\begin{milestonebox}[Files and directories]
\begin{checklist}
  \item The VFS resolves paths across \code{/} (tar), \code{/dev} (devfs) and \code{/mnt} (TurkFS).
  \item Every process has descriptors 0--2 on \code{/dev/tty}; \code{fork}, \code{dup2} and \code{exec} handle descriptors correctly.
  \item Programs are loaded by path from the file system.
  \item TurkFS reads and writes files and directories, and \code{fsck.minix -f} on the host finds no errors.
  \item All file-system tests pass, including the 2\,MiB file.
\end{checklist}
\end{milestonebox}

\begin{questionbox}
\begin{enumerate}
  \item List the disk reads needed to open and read \code{/home/turk/notes.txt} with an empty cache, following \OSTEP\,\S40.4.
  \item What is the largest file TurkFS can hold with seven direct zones, one single and one double indirect block, 1\,KiB zones and 4-byte zone numbers?
  \item What is the difference between a file descriptor, an open file and an inode? Which of them does \code{fork} copy?
  \item What does the VFS layer buy you? Name three things in Turk-OS that would be harder without it.
  \item Why is crash consistency hard? In what order should \code{create} write its blocks, and why?
  \item What is the difference between a hard link and a symbolic link? How does the link count relate to \code{unlink}?
\end{enumerate}
\end{questionbox}

\section{Going further}

Make TurkFS the root file system and keep the tar archive only for early boot. Add a simple metadata journal
(\OSTEP\,\S42.3): write the blocks of each operation to a log region first, then a commit block, then apply them, and
replay the log at mount. Write a FAT16 driver (\MOS\,\S4.5.1) so Turk-OS can read USB-stick images made on any
computer. Or add symbolic links, \code{rename}, and file permissions checked against the process's user ID, which
Phase 15 builds on.

\nextphase{14}{Userland: libc, Shell and Utilities}
\booklegend
\end{document}
